% ===================================================================
%  도표 제 12 호 — 정거장. 철길. 배.
%
%  Traduction, faite en 2026, en coreen d'aujourd'hui, sur l'ido de
%  texto/io. Regles de la colonne : voir 10-dopyo-01.tex. Une seule
%  suite d'alineas ; le « Noto. » d'ouverture appartient au bloc apar,
%  et la note de bas de page du feuillet 85 a son bloc noto.
%
%  LE RENVOI (*) N'EST PAS UN NUMERO ET SE PORTE QUAND MEME. Le bloc
%  19 appelle la note par « \textsuperscript{(*)} » ; renvoji.py le
%  releve comme les autres, et la colonne l'ecrit tel quel.
%
%  LE COREEN DU RAIL EST ENTIER, ET IL EST PRESQUE AUSSI VIEUX QUE LE
%  LIVRET. Le premier train coreen roule en 1899, vingt-sept ans avant
%  ce texte : la langue a eu le temps de se faire ses mots. 철길 la
%  voie, 객차 la voiture, 기적 le sifflet, 굴 le tunnel, 갈림길
%  l'aiguillage, 짐꾼 le porteur, 시간표 l'indicateur, 침대차 le
%  wagon-lit, 식당차 le wagon-restaurant, 우편차 le wagon-poste,
%  화물차 le wagon de marchandises, 차장 le chef de train, 개찰구 le
%  guichet. Restent en anglais les objets que le coreen d'aujourd'hui
%  nomme en anglais : 플랫폼, 보일러, 피스톤, 트렁크.
%
%  C'est le troisieme tableau de suite ou rien ne manque, apres le
%  port et la ville. La coupure du vocabulaire, dans cette colonne, ne
%  passe jamais par la langue : elle passe par la date d'entree de
%  l'objet dans le pays, exactement comme au tableau 13 telougou.
%
%  TRENTE PAIRES. Le tableau raconte un voyage et enumere ce qui
%  defile par la fenetre ; il attache surtout dans les descriptions de
%  quai, ou chaque objet est pose par rapport a un autre.
% ===================================================================

%%K t12-apar-1 apar
\VUsaut{4.0mm}
\VUtitre{15.0pt}{60}{도표 제 12 호}
\VUsaut{2.4mm}
\VUpk{11.4pt}{40}{정거장. --- 철길. --- 배.}
\VUsaut{3.4mm}
\textsc{알림.} --- \textit{나라마다 시각을 세는 법이 다르므로, 우리는 보기로}
\VUgras{프랑스 도표} \textit{의 시각을 쓴다.}

%%K t12-01-1 p
1. --- 나는 방금 독일을 거쳐 여행했다. 떠나기 전에 기차가 몇 시에 떠나는지 알려고
\VUgras{시간표 책} \textsuperscript{(15)}을 샀다. 가장 알맞은 기차가 파리에서 밤 아홉
시에 떠나 스트라스부르에 한 시 십삼 분에 닿는다는 것을 확인한 뒤, 나는 여행 채비를
했고, 이튿날 정거장으로 데려다 달라고 했다.

%%K t12-02-1 p
2. --- 큰 출발 대합실로 들어섰을 때, 시골 \VUgras{신부님} \textsuperscript{(2)} 뒤로 —
그분은 손에 제 \VUgras{여행 가방} \textsuperscript{(3)}을 들고 계셨다 — 이미 많은
\VUgras{나그네} \textsuperscript{(1)}가 와 있었다.

%%K t12-02-2 p
나는 \VUgras{전보} \textsuperscript{(5)}를 부치고 싶었다. 그래서 \VUgras{검표원}
\textsuperscript{(6)}에게 물어 \VUgras{전신국} \textsuperscript{(4)}이 어디인지 알아
두었다.

%%K t12-03-1 p
3. --- \VUgras{대합실} \textsuperscript{(7)}로 들어가기 전에 나는 돌아오는 길까지 넣은
\VUgras{표} \textsuperscript{(9)}를 끊었다.

%%K t12-04-1 p
4. --- \VUgras{개찰구} \textsuperscript{(8)} 앞에서 오래 기다리지 않았다. 사람이 적었기
때문이다. 휴가를 떠나는 \VUgras{병사} \textsuperscript{(10)} 하나가 친절히 제 차례를 나에게
양보해 주었다. 덕분에 나는 \VUgras{역장} \textsuperscript{(13)}에게 몇 가지 물어볼 짬이
있었다 — 그분은 감시 \VUgras{경감} \textsuperscript{(12)}과, 당직 \VUgras{헌병}
\textsuperscript{(11)}과 이야기를 나누고 계셨다.

%%K t12-tit-1 sub
\VUpk{10.2pt}{40}{짐 부치기.}
\VUsaut{3.0mm}

%%K t12-05-1 p
5. --- \VUgras{책 가게} \textsuperscript{(14)}에서 신문을 사고 나서, 나는 내 \VUgras{짐}
\textsuperscript{(17)}의 무게를 달게 했다. 그 짐은 \VUgras{짐꾼} \textsuperscript{(16)}이
방금 \VUgras{짐수레} \textsuperscript{(19)}로 날라 온 것이다. \VUgras{서기}
\textsuperscript{(22)}가 — \VUgras{저울} \textsuperscript{(21)}을 맡은 사람이다 — 내 궤가
서른다섯 킬로그램이라고 알려 주었다. 그러니 다섯 킬로그램의 \VUgras{초과 무게}가 되고,
그 값을 \VUgras{짐 창구} \textsuperscript{(23)}에서 따로 물어야 한다.

%%K t12-06-1 p
6. --- 너무 일찍 온 터라 정거장에서 나그네들이 오가는 것을 살펴볼 짬이 있었다. 어떤
이들은 대수롭지 않은 짐을 \VUgras{짐 맡기는 곳} \textsuperscript{(25)}에 맡기거나 찾아
갔다. 다른 이들은 \VUgras{시간표} \textsuperscript{(26)}를 들여다보았다. 막 도착한
나그네들은 \VUgras{나가는 곳} \textsuperscript{(32)}으로 서둘러 갔다 — 손에
\VUgras{트렁크} \textsuperscript{(24)}를 들고. 몇몇은 \VUgras{짐 내주는 곳}
\textsuperscript{(27)}에서 기다리며 제 \VUgras{쪽지} \textsuperscript{(29)}를 내밀었다 —
제 궤와 \VUgras{상자} \textsuperscript{(28)}와 \VUgras{바구니} \textsuperscript{(30)}를
찾으려는 것이다.

%%K t12-07-1 p
7. --- \VUgras{세관 관리} \textsuperscript{(33)}들이 꾸러미를 꼼꼼히 살폈다. 신고할 것이
있는지 확인하려는 것이다.

%%K t12-08-1 p
8. --- 어떤 나그네들은 \VUgras{간이식당} \textsuperscript{(34)}으로 갔다. 거기에서
\VUgras{종업원} \textsuperscript{(35)}이 그들에게 부지런히 시중을 든다. 그들은 서둘러야
한다. \VUgras{기차} \textsuperscript{(36)}가 급행이라 정거장에 오래 서 있지 않기 때문이다.

%%K t12-09-1 p
9. --- 그 다음에 나는 출발 승강장으로 가서 곧바로 가는 \VUgras{객차} \textsuperscript{(37)}
를 찾았다. 여행 도중에 차를 갈아타지 않으려는 것이다. 나는 복도가 있는 차에 올랐는데,
그 안에는 담배 피우는 이들을 위한 \VUgras{칸} \textsuperscript{(38)}과 「혼자 오신
부인들」을 위해 남겨 둔 칸이 있다.

%%K t12-tit-2 sub
\VUpk{10.2pt}{40}{밤에 떠나기.}
\VUsaut{3.0mm}

%%K t12-10-1 p
10. --- \VUgras{디딤판} \textsuperscript{(40)}을 오르고 \VUgras{문} \textsuperscript{(39)}을
닫고 \VUgras{차창 가리개} \textsuperscript{(41)}를 말아 올린 뒤, 작은 짐을 \VUgras{그물 선반}
\textsuperscript{(43)} 위에 얹고 \VUgras{긴 의자} \textsuperscript{(42)}에 앉았다.
\VUgras{등 켜는 이} \textsuperscript{(44)}가 등을 켜는 것을 보고, 밤을 객차에서 보내야
한다는 생각이 들어 \VUgras{베개} \textsuperscript{(45)} 빌려주는 일을 맡은 직원을 불러
하나를 달라고 했다.

%%K t12-11-1 p
11. --- 표를 검사하러 검표원이 왔다. 시각이 되었는데 왜 아직 떠나지 않느냐고 물었다.
\VUgras{차장} \textsuperscript{(46)}이 \VUgras{수하물차} \textsuperscript{(47)}에 자전거를
싣게 하느라 바쁘다고 했다. 게다가 \VUgras{발 데우개} \textsuperscript{(48)}를 아직 다
갈지 못했다고 했다.

%%K t12-12-1 p
12. --- 그러나 우리는 곧 떠날 참이었다. \VUgras{화부} \textsuperscript{(50)}가 제
\VUgras{석탄차} \textsuperscript{(49)} 위에서 석탄을 집어 \VUgras{기관차} \textsuperscript{(54)}
의 불을 돋우었고, \VUgras{기관사} \textsuperscript{(51)}는 \VUgras{부역장} \textsuperscript{(52)}
의 신호만 기다리고 있었기 때문이다. 그 부역장은 \VUgras{승강장} \textsuperscript{(53)} 위에
있었다.

%%K t12-13-1 p
13. --- 기관차의 \VUgras{보일러} \textsuperscript{(56)}가 무겁게 울렸다. \VUgras{굴뚝}
\textsuperscript{(55)}이 \VUgras{김} \textsuperscript{(60)}과 섞인 연기를 쏟아 냈다.
날카로운 \VUgras{기적} \textsuperscript{(59)}이 크게 울렸고, \VUgras{피스톤}
\textsuperscript{(58)}이 움직이기 시작하여 그 무거운 기계의 \VUgras{바퀴}
\textsuperscript{(57)}를 밀었다.

%%K t12-tit-3 sub
\VUsaut{3.0mm}
\VUpk{10.2pt}{40}{출발!}
\VUsaut{3.0mm}

%%K t12-14-1 p
14. --- \VUgras{철길} \textsuperscript{(64)} 위를 달리는 기차의 빠르기가 더해졌다.
\VUgras{승강장} \textsuperscript{(65)}이 사라졌다. \VUgras{뒷간} \textsuperscript{(66)}을
지나쳤고, 건너편 \VUgras{선로} \textsuperscript{(63)}에 세워 둔 차들도 지나쳤다.
\VUgras{침대차} \textsuperscript{(67)}, \VUgras{식당차} \textsuperscript{(68)},
\VUgras{우편차} \textsuperscript{(69)}였다.

%%K t12-noto-1 noto
\VUnotes{143.2mm}{%
\textit{(*) 알림.} --- \textit{이 여행 이야기는 전쟁 이전의 국경에 따른 것임을
알아야 한다.}%
}

%%K t12-15-1 p
15. --- 그 다음에 \VUgras{화물 정거장} \textsuperscript{(71)}이 우리 눈앞을 지나갔다.
지붕 아래에서 직원들이 완행으로 부친 \VUgras{화물} \textsuperscript{(72)}을 싣고 있었다.
\VUgras{작업반} \textsuperscript{(76)}이 — \VUgras{물탱크} \textsuperscript{(73)} 가까이에서
— \VUgras{가축차} \textsuperscript{(75)}와 \VUgras{무개차} \textsuperscript{(79)}를 움직여
\VUgras{화물 열차} \textsuperscript{(80)}를 꾸몄다.

%%K t12-16-1 p
16. --- 우리는 마지막 \VUgras{갈림길} \textsuperscript{(78)}을 지났다. 그 곁에 전철수가
서 있었다. \VUgras{신호판} \textsuperscript{(86)}을 지나치자 정거장이 멀리 사라졌다.

%%K t12-17-1 p
17. --- 곧 우리는 \VUgras{쇠다리} \textsuperscript{(74)}를 건넜다. 그것이 \VUgras{강}
\textsuperscript{(81)}을 가로지른다. 그 다리를 건너고 나서 우리는 강을 따라갔는데, 그
위로 힘센 \VUgras{예인선} \textsuperscript{(82)}들이 묵직한 \VUgras{짐배} \textsuperscript{(83)}
를 끌고 있었다. 또 \VUgras{여객선} \textsuperscript{(84)} 한 척도 보았다 — 그것이
\VUgras{뜬 부두} \textsuperscript{(85)}에 닿으려는 참이었다.

%%K t12-18-1 p
18. --- 여러 정거장을 아주 빠르게 지나친 뒤 우리는 몇 개의 \VUgras{굴} \textsuperscript{(87)}
을 지났고, \VUgras{건널목 지킴이}들의 헤아릴 수 없는 \VUgras{집} \textsuperscript{(88)}
들을 뒤에 남겼다. 기차의 흔들림에 요람처럼 흔들리며 나는 깊이 잠들었다.

%%K t12-tit-4 sub
\VUsaut{3.0mm}
\VUpk{10.2pt}{40}{도이치아브리쿠르 정거장.}
\VUsaut{3.0mm}

%%K t12-19-1 p
19. --- 외치는 직원의 목소리에 나는 갑자기 깨어났다. 「도이치아브리쿠르!
\textsuperscript{(*)} 모두 내리시오!」 세관 검사와 국경의 귀찮은 절차가 다 있었다. 나는
내 회중시계를 정거장 \VUgras{큰 시계} \textsuperscript{(89)}에 맞추어야 했다. 그 시계가
쉰다섯 분 빨랐다. 잠이 덜 깬 나는 \VUgras{긴 바늘} \textsuperscript{(90)}과
\VUgras{짧은 바늘} \textsuperscript{(91)}을 겨우 가려보았다. \VUgras{글자판}
\textsuperscript{(92)} 자체가 아침 안개로 온통 흐릿했다.

%%K t12-20-1 p
20. --- 세관 사람들이 검사하는 동안, 나는 하품하며 정거장 벽을 꾸미는 \VUgras{벽보}
\textsuperscript{(93)}를 읽었다. 시간을 보내려고 아침을 먹고, 스트라스부르가 아직 얼마나
먼지 보려고 독일 \VUgras{철도망} \textsuperscript{(94)} 지도를 들여다보았다. 그리고 여행의
끝에 곧 이르리라는 바람으로 힘을 얻어 다시 내 객차로 올라갔다.

\VUsaut{6.0mm}
\VUfilet{22mm}
